\chapter{Klaus Friedrich Roth --- Fields Medal 1958}\label{ch1958klausfriedrichroth:klaus-friedrich-roth-fields-medal-1958-a-comprehensive-report}

\begin{center}\rule{0.5\linewidth}{0.5pt}\end{center}

\textbf{Prepared:} July 2026 \textbf{Subject:} Klaus Friedrich Roth (29
October 1925 -- 10 November 2015) \textbf{Position:} Professor,
University College London, then Imperial College London \textbf{Doctoral
Advisor:} Theodor Estermann (University College London; informal mentor:
Harold Davenport) \textbf{Fields Medal Citation:} \emph{``Solved in 1955
the famous Thue--Siegel problem concerning the approximation to
algebraic numbers by rational numbers and proved in 1952 that a sequence
with no three numbers in arithmetic progression has zero density (a
conjecture of Erdős and Turán of 1935).''}

\begin{center}\rule{0.5\linewidth}{0.5pt}\end{center}

\section{1. Introduction}\label{ch1958klausfriedrichroth:introduction}

Klaus Friedrich Roth was the first British mathematician to win the
Fields Medal, awarded at the 1958 International Congress of
Mathematicians in Edinburgh. He is remembered for two monumental
achievements in number theory: \textbf{Roth's theorem on Diophantine
approximation}, which settled a century-old problem on approximating
algebraic numbers by rationals, and his proof that sequences with no
three-term arithmetic progression have zero density, a foundational
result in additive combinatorics.

\begin{prereqs} Elementary number theory, basic real analysis (sequences, convergence), pigeonhole principle. \end{prereqs}

\begin{objectives} State Roth's theorem on diophantine approximation; explain why it is a deep generalization of Liouville's theorem; describe the application to the Thue-Siegel-Roth conjecture. \end{objectives}


Beyond these well-known results, Roth made deep contributions to: - The
\textbf{large sieve} and its applications to primes in arithmetic
progressions. - The \textbf{theory of discrepancy} (irregularities of
distribution), achieving the optimal lower bound for the discrepancy of
sequences in the unit square. - The \textbf{theory of integer sequences}
and quantitative versions of Szemerédi-type theorems. -
\textbf{Diophantine approximation} and the geometry of numbers.

This report provides a comprehensive overview of Roth's life, his
monumental contributions, and the lasting impact of his work. Roth was
primarily a problem-solver, described by Davenport as demonstrating that
``the great unsolved problems of mathematics may still yield to direct
attack, however difficult and forbidding they appear.''

\begin{center}\rule{0.5\linewidth}{0.5pt}\end{center}

\section{2. Biographical Background}\label{ch1958klausfriedrichroth:biographical-background}

\subsection{2.1 Early Life and
Education}\label{ch1958klausfriedrichroth:early-life-and-education}

Klaus Friedrich Roth was born on 29 October 1925 in Breslau, Lower
Silesia, Prussia (now Wrocław, Poland), to Franz Roth, a solicitor, and
Matilde Roth (née Liebrecht). To escape Nazi persecution, the family
fled Germany in 1933 and settled in London. Franz Roth had suffered gas
poisoning during World War I and died a few years after their arrival in
England.

Roth attended \textbf{St Paul's School} between 1937 and 1943, where the
school was evacuated to Easthampstead Park in Berkshire during the war.
There he excelled in mathematics and chess. His mathematics master, Mr
Dowswell, observed that Roth possessed ``complete intellectual
honesty.''

He went up to \textbf{Peterhouse, Cambridge} in 1943 to read
mathematics. However, unbearable nervousness severely hampered his
examination performance, and he graduated with third-class honours in
1945 --- a poor start for a future Fields medallist.

\subsection{2.2 War Service and Early
Career}\label{ch1958klausfriedrichroth:war-service-and-early-career}

After Cambridge, Roth undertook war service as an alien and became a
junior master at \textbf{Gordonstoun School} in Scotland. He divided his
leisure time between riding a powerful motorcycle through the Scottish
countryside and playing chess with future British Chess Champion Robert
Combe.

Peterhouse did not support Roth's return to Cambridge, and his tutor
John Charles Burkill suggested he pursue a commercial job. However,
\textbf{Harold Davenport} recognized Roth's exceptional mathematical
talent and arranged for him to study at \textbf{University College
London} (UCL).

\subsection{2.3 Academic Career}\label{ch1958klausfriedrichroth:academic-career}

\begin{itemize}
\tightlist
\item
 \textbf{1946--1948:} Undergraduate studies at UCL, receiving a
 master's degree in 1948.
\item
 \textbf{1948--1950:} Assistant lecturer, UCL.
\item
 \textbf{1950--1961:} Lecturer, UCL.
\item
 \textbf{1955:} Published Roth's theorem on Diophantine approximation.
\item
 \textbf{1958:} Fields Medal, ICM Edinburgh.
\item
 \textbf{1960:} Elected Fellow of the Royal Society.
\item
 \textbf{1961--1966:} Reader, then Professor, UCL.
\item
 \textbf{1966--1988:} Professor of Pure Mathematics, Imperial College
 London.
\item
 \textbf{1988:} Emeritus Professor, Imperial College.
\item
 \textbf{2005--2015:} Honorary Fellow, Peterhouse, Cambridge.
\end{itemize}

\subsection{2.4 Personality and
Approach}\label{ch1958klausfriedrichroth:personality-and-approach}

Roth was known for his modesty, intellectual honesty, and intense focus
on solving hard problems. He did not build elaborate theories but rather
attacked central unsolved problems with tremendous technical ingenuity.
His work was characterized by startlingly original ideas that often
opened entirely new fields of research.

\begin{center}\rule{0.5\linewidth}{0.5pt}\end{center}

\section{3. Roth's Theorem on Diophantine
Approximation}\label{ch1958klausfriedrichroth:roths-theorem-on-diophantine-approximation}

\subsection{3.1 Historical Background}\label{ch1958klausfriedrichroth:historical-background}

The problem of approximating irrational numbers by rationals has a long
history. For a real number \(\alpha\), the \textbf{approximation
exponent} \(\mu(\alpha)\) is the supremum of \(\mu\) such that:

\[\left|\alpha - \frac{p}{q}\right| < \frac{1}{q^\mu}\]

has infinitely many rational solutions \(p/q\).

\begin{itemize}
\tightlist
\item
 \textbf{Dirichlet (1842):} For any irrational \(\alpha\), there are
 infinitely many \(p/q\) with \(|\alpha - p/q| < 1/q^2\), so
 \(\mu(\alpha) \geq 2\).
\item
 \textbf{Liouville (1844):} For algebraic numbers of degree \(d\),
 \(\mu(\alpha) \leq d\), giving the first proof of transcendental
 numbers.
\item
 \textbf{Thue (1909):} For algebraic \(\alpha\) of degree \(d \geq 3\),
 \(\mu(\alpha) \leq d/2 + 1\).
\item
 \textbf{Siegel (1921):} \(\mu(\alpha) \leq 2\sqrt{d}\).
\item
 \textbf{Dyson (1947):} \(\mu(\alpha) \leq \sqrt{2d}\).
\item
 \textbf{Gelfond (1952):} Independent proof of the Dyson bound using
 Pell-type equations.
\end{itemize}

Each improvement reduced the exponent but still left a dependence on the
degree \(d\). Siegel conjectured that the true exponent was \(2\) for
all algebraic numbers.

\subsection{3.2 Statement}\label{ch1958klausfriedrichroth:statement}

\textbf{Roth's Theorem (1955):} For any algebraic irrational number
\(\alpha\) and any \(\epsilon > 0\), there are only finitely many
rational numbers \(p/q\) such that:

\[\left|\alpha - \frac{p}{q}\right| < \frac{1}{q^{2+\epsilon}}\]

Equivalently, the approximation exponent \(\mu(\alpha) = 2\) for every
algebraic irrational \(\alpha\).

\subsection{3.3 The Proof}\label{ch1958klausfriedrichroth:the-proof}

Roth's proof combined several deep ideas:

\begin{enumerate}
\def\labelenumi{\arabic{enumi}.}
\item
 \textbf{A gap principle:} Construct a sequence of good approximations
 \(p_1/q_1, p_2/q_2, \ldots\) and show that the denominators must grow
 extremely rapidly.
\item
 \textbf{The method of hyperplanes (Roth's lemma):} A combinatorial
 result about the vanishing of homogeneous polynomials at many rational
 points. Roth proved that if \(P(x_1, \ldots, x_n)\) is a polynomial
 with integer coefficients and bounded degree, and if \(P\) vanishes at
 many rational points \((a_1/b, \ldots, a_n/b)\), then the polynomial
 must have a special factorization.
\item
 \textbf{Construction of an auxiliary polynomial:} This polynomial,
 built from the sequence of approximations, vanishes to high order at
 many points, leading to a contradiction unless the approximations are
 bounded in number.
\end{enumerate}

\subsection{3.4 The Role of Ridout}\label{ch1958klausfriedrichroth:the-role-of-ridout}

\textbf{David Ridout} later extended Roth's theorem to allow \(p\) and
\(q\) to have restricted prime factorizations, establishing the
\textbf{Roth--Ridout theorem}. This refinement has important
applications in transcendence theory.

\subsection{3.5 The Three-Step Strategy}\label{ch1958klausfriedrichroth:the-three-step-strategy}

Roth's proof follows a three-step strategy that is now standard in
Diophantine approximation:

\begin{enumerate}
\def\labelenumi{\arabic{enumi}.}
\item
 \textbf{Construction of an auxiliary polynomial:} Given good
 approximations \(p_1/q_1, \ldots, p_m/q_m\), construct a polynomial
 \(P(x_1, \ldots, x_m)\) with integer coefficients that vanishes to
 high order at \((p_1/q, \ldots, p_m/q)\) for a common denominator
 \(q\).
\item
 \textbf{Application of Roth's lemma:} Show that \(P\) cannot vanish at
 too many such points unless it has a special factorization, which is
 impossible because of size constraints.
\item
 \textbf{Contradiction via size estimates:} Use the fact that \(P\) has
 bounded coefficients and the \(p_i/q_i\) are good approximations to
 obtain a contradiction unless \(m\) is bounded.
\end{enumerate}

\subsection{3.6 Significance}\label{ch1958klausfriedrichroth:significance}

Roth's theorem is a landmark in number theory: - It completely solved
the Thue--Siegel problem, which had resisted the efforts of the greatest
number theorists for half a century. - It shows that algebraic numbers
are as poorly approximable by rationals as almost all real numbers (in
the sense of Lebesgue measure). - It is a key ingredient in the theory
of \textbf{transcendental numbers} and \textbf{Diophantine geometry}. -
Siegel wrote that the result ``will be remembered as long as mankind is
interested in mathematics.'' - Davenport, in his Fields Medal speech,
said: ``The achievement is one that speaks for itself: it closes a
chapter, and a new chapter is opened.''

\subsection{3.7 The Thue--Siegel--Roth
Method}\label{ch1958klausfriedrichroth:the-thuesiegelroth-method}

Roth's method completed a line of research known as the
\textbf{Thue--Siegel--Roth method}: - \textbf{Thue (1909):} First
non-trivial bound using an auxiliary polynomial in one variable. -
\textbf{Siegel (1921):} Introduction of the \textbf{Siegel lemma}
(existence of polynomials with small integer coefficients). -
\textbf{Dyson (1947):} Use of multiple variables in the auxiliary
polynomial. - \textbf{Roth (1955):} Optimal exponent 2 using a large
number of variables and Roth's lemma.

\begin{center}\rule{0.5\linewidth}{0.5pt}\end{center}

\section{4. Sequences Without Three-Term Arithmetic
Progressions}\label{ch1958klausfriedrichroth:sequences-without-three-term-arithmetic-progressions}

\subsection{4.1 The Erdős--Turán
Conjecture}\label{ch1958klausfriedrichroth:the-erdux151sturuxe1n-conjecture}

In 1936, Erdős and Turán conjectured that any subset of the natural
numbers with positive upper density contains arbitrarily long arithmetic
progressions. As a first step, they asked whether a sequence with no
three-term arithmetic progression must have zero density.

\subsection{4.2 Roth's Theorem}\label{ch1958klausfriedrichroth:roths-theorem}

In 1952, Roth proved:

\begin{quote}
\textbf{Roth's Theorem (on arithmetic progressions):} For any set
\(A \subset \{1, \ldots, N\}\) with no three distinct elements
\(x, y, z\) satisfying \(x + z = 2y\), we have:
\end{quote}

\[|A| = O\left(\frac{N}{\log \log N}\right)\]

In particular, the upper density of a set with no three-term arithmetic
progression is zero. This was the first significant progress on the
Erdős--Turán conjecture.

\subsection{4.3 The Method}\label{ch1958klausfriedrichroth:the-method}

Roth's proof used \textbf{Fourier analysis} (the Hardy--Littlewood
circle method) combined with a density increment argument:

\begin{enumerate}
\def\labelenumi{\arabic{enumi}.}
\tightlist
\item
 Suppose \(A \subset \{1, \ldots, N\}\) has density \(\delta = |A|/N\)
 and no 3-term arithmetic progression.
\item
 Compute the Fourier coefficients
 \(\widehat{1_A}(r) = \sum_{n \in A} e(nr/N)\).
\item
 If all non-trivial Fourier coefficients are small, then by a
 \textbf{van der Corput} estimate, \(A\) contains many 3-term
 progressions.
\item
 If some Fourier coefficient is large, then \(A\) has increased density
 on a long arithmetic progression.
\item
 Iterating gives the bound \(|A| = O(N/\log \log N)\).
\end{enumerate}

\subsection{4.4 Later Refinements}\label{ch1958klausfriedrichroth:later-refinements}

Roth's bound was improved over decades: - \textbf{Szemerédi (1975):}
\(|A| = o(N)\) for \(k\)-term progressions (Szemerédi's theorem). -
\textbf{Heath-Brown and Szemerédi:} \(|A| = O(N/(\log N)^c)\). -
\textbf{Bourgain (2008):} \(|A| = O(N(\log \log N)^2/(\log N)^{2/3})\).
- \textbf{Bloom and Sisask (2020):} \(|A| = O(N/(\log N)^{1+c})\) for
some \(c > 0\).

Roth's Fourier-analytic method remains the foundation for the study of
three-term arithmetic progressions.

\begin{center}\rule{0.5\linewidth}{0.5pt}\end{center}

\section{5. The Large Sieve}\label{ch1958klausfriedrichroth:the-large-sieve}

\subsection{5.1 Background}\label{ch1958klausfriedrichroth:background}

The \textbf{large sieve} is a method in analytic number theory for
bounding the size of a set of integers that avoid certain residue
classes modulo primes. Its early development was due to Linnik, but
Roth's 1965 paper was a seminal contribution.

\subsection{5.2 Roth's Contribution}\label{ch1958klausfriedrichroth:roths-contribution}

Roth's paper ``On the large sieves of Linnik and Rényi'' developed a
quantitative version of the large sieve with applications to the
distribution of primes in arithmetic progressions. He proved:

\begin{quote}
\textbf{Roth's Large Sieve Inequality:} For a set
\(A \subset \{1, \ldots, N\}\) and primes \(p \leq Q\), let \(Z(p)\) be
the number of residue classes modulo \(p\) that contain no element of
\(A\). Then:
\end{quote}

\[|A| \leq \frac{N + Q^2}{\sum_{p \leq Q} Z(p)/p}\]

\subsection{5.3 The Linnik--Rényi--Roth
Inequality}\label{ch1958klausfriedrichroth:the-linnikruxe9nyiroth-inequality}

Roth's version of the large sieve improved on earlier work of
\textbf{Linnik} and \textbf{Rényi}. The key innovation was a
\textbf{combinatorial inequality} that avoids the analytic number theory
used by previous authors. Roth's inequality is:

\[\sum_{q \leq Q} \mu(q)^2 \sum_{a=1, (a,q)=1}^q \left|\sum_{n \leq N} a_n e(an/q)\right|^2 \leq (N + Q^2) \sum_{n \leq N} |a_n|^2\]

\subsection{5.4 Applications}\label{ch1958klausfriedrichroth:applications}

The large sieve is essential for: - The \textbf{Bombieri--Vinogradov
theorem} on primes in arithmetic progressions. - \textbf{Barban's
theorem} and related results on the mean square of the error term. -
\textbf{The Brun--Titchmarsh theorem} and its generalizations. -
\textbf{Sieving for prime tuples} and the Hardy--Littlewood conjectures.
- The \textbf{Gallagher prime number theorem} and the
\textbf{Barban--Davenport--Halberstam theorem}.

\begin{center}\rule{0.5\linewidth}{0.5pt}\end{center}

\section{6. The Theory of Discrepancy}\label{ch1958klausfriedrichroth:the-theory-of-discrepancy}

\subsection{6.1 The Problem}\label{ch1958klausfriedrichroth:the-problem}

The \textbf{discrepancy} of a sequence \(x_1, \ldots, x_N \in [0,1]\)
measures how far its empirical distribution is from uniform:

\[D_N = \sup_{0 \leq a < b \leq 1} \left|\frac{\#\{n \leq N : x_n \in [a,b]\}}{N} - (b-a)\right|\]

van der Corput (1935) conjectured that \(D_N \to \infty\) for any
sequence, i.e., perfect equidistribution is impossible.

\subsection{6.2 Roth's Lower Bound}\label{ch1958klausfriedrichroth:roths-lower-bound}

In 1954, Roth proved the optimal lower bound:

\begin{quote}
\textbf{Roth's Discrepancy Theorem:} For any sequence
\(x_1, \ldots, x_N \in [0,1]\), the discrepancy satisfies:
\end{quote}

\[D_N \geq c \frac{\log N}{N}\]

for some absolute constant \(c > 0\). Moreover, this bound is sharp:
there exist sequences with \(D_N = O(\log N/N)\).

\subsection{6.3 The Method}\label{ch1958klausfriedrichroth:the-method-1}

Roth reformulated the discrepancy problem in \textbf{geometric terms}.
He considered the \textbf{discrepancy function} \(f(x,y)\) on the unit
square and proved a lower bound on its \(L^2\) norm using
\textbf{Fourier analysis} and \textbf{orthogonal functions}. The key
estimate:

\[\int_0^1 \int_0^1 |f(x,y)|^2 \, dx \, dy \geq c \frac{(\log N)^2}{N^2}\]

\subsection{6.4 The Geometric
Reformulation}\label{ch1958klausfriedrichroth:the-geometric-reformulation}

Roth's key insight was to reformulate discrepancy in terms of a
\textbf{Hilbert space} problem. Let \(\chi_R\) be the characteristic
function of a rectangle \(R = [0,a) \times [0,b)\), and let \(f_N\) be
the discrepancy function:

\[f_N(a,b) = \frac{1}{N}\sum_{n=1}^N \chi_R(x_n,y_n) - ab\]

Then the \(L^2\) norm of \(f_N\) over \([0,1]^2\) is bounded below by
\(c(\log N)/N\). The proof uses the \textbf{orthogonal Haar system} or
\textbf{Walsh functions} to decompose the rectangle indicator functions
and exploit their orthogonality.

\subsection{6.5 Higher Dimensions}\label{ch1958klausfriedrichroth:higher-dimensions}

Roth's method extends to higher dimensions. For \(d\)-dimensional
sequences:

\[D_N \geq c_d \frac{(\log N)^{(d-1)/2}}{N}\]

These bounds are optimal up to constant factors, and have applications
in \textbf{numerical integration} (quasi-Monte Carlo methods) and
\textbf{computational finance}. The optimality was established by
constructions of \textbf{Halton}, \textbf{Hammersley}, and
\textbf{Faure}.

\begin{center}\rule{0.5\linewidth}{0.5pt}\end{center}

\section{7. Roth's Theorem on Rational Approximations of Algebraic
Numbers}\label{ch1958klausfriedrichroth:roths-theorem-on-rational-approximations-of-algebraic-numbers}

\subsection{7.1 Effective vs.~Ineffective
Bounds}\label{ch1958klausfriedrichroth:effective-vs.-ineffective-bounds}

Roth's theorem is \textbf{ineffective}: it proves finiteness but does
not give a computable bound on the size of the largest approximating
denominator. This is a fundamental limitation of the method.

\subsection{7.2 Effective Improvements}\label{ch1958klausfriedrichroth:effective-improvements}

\begin{itemize}
\tightlist
\item
 \textbf{Baker's theorem} (1968, Alan Baker): Gives effective bounds
 for linear forms in logarithms, which can be used to bound the largest
 solution to Thue equations.
\item
 \textbf{Feldman (1971):} For \(|\alpha - p/q| < q^{-\kappa}\) with
 \(\kappa < 2\), there is an effective bound depending on \(\kappa\)
 and the height of \(\alpha\).
\end{itemize}

\subsection{7.3 The Subspace Theorem}\label{ch1958klausfriedrichroth:the-subspace-theorem}

\textbf{Schmidt's subspace theorem} (1972) generalizes Roth's theorem to
simultaneous approximation: for linearly independent linear forms
\(L_1, \ldots, L_n\) with algebraic coefficients and any
\(\epsilon > 0\), the inequality:

\[\prod_{i=1}^n |L_i(\mathbf{x})| < \|\mathbf{x}\|^{-\epsilon}\]

has only finitely many solutions \(\mathbf{x} \in \mathbb{Z}^n\) outside
a finite union of subspaces. This powerful generalization has many
applications in Diophantine geometry.

\begin{center}\rule{0.5\linewidth}{0.5pt}\end{center}

\section{8. Roth's Work on the Sieve of
Eratosthenes}\label{ch1958klausfriedrichroth:roths-work-on-the-sieve-of-eratosthenes}

\subsection{8.1 A Refined Sieve}\label{ch1958klausfriedrichroth:a-refined-sieve}

Roth obtained refined estimates for the \textbf{sieve of Eratosthenes},
improving bounds on the number of integers up to \(x\) that are free of
small prime factors. These estimates are crucial for understanding the
distribution of prime numbers.

\subsection{8.2 Applications}\label{ch1958klausfriedrichroth:applications-1}

\begin{itemize}
\tightlist
\item
 \textbf{The Brun--Titchmarsh theorem} gives upper bounds for
 \(\pi(x; q, a)\), the number of primes \(p \leq x\) in the arithmetic
 progression \(p \equiv a \pmod{q}\).
\item
 Roth's sieve estimates improved the constants in these bounds.
\item
 The \textbf{large sieve} and its variants (Selberg sieve, Bombieri
 sieve) all benefit from Roth's foundational work.
\end{itemize}

\begin{center}\rule{0.5\linewidth}{0.5pt}\end{center}

\section{9. Diophantine Approximation and the Geometry of
Numbers}\label{ch1958klausfriedrichroth:diophantine-approximation-and-the-geometry-of-numbers}

\subsection{9.1 Minkowski's Theorem}\label{ch1958klausfriedrichroth:minkowskis-theorem}

Roth made contributions to the \textbf{geometry of numbers}, a subject
pioneered by Minkowski. The central result is Minkowski's convex body
theorem: a symmetric convex body of volume \(> 2^n\) in \(\mathbb{R}^n\)
contains a non-zero integer lattice point.

\subsection{9.2 Roth's Theorem on Linear
Forms}\label{ch1958klausfriedrichroth:roths-theorem-on-linear-forms}

Roth proved a theorem on the simultaneous approximation of linear forms,
extending the results of Dirichlet and Kronecker. For a system of \(m\)
linear forms in \(n\) variables with real coefficients, Roth's
inequality gives:

\[\min_{\mathbf{x} \in \mathbb{Z}^n \setminus \{0\}} \max_{1 \leq i \leq m} \|L_i(\mathbf{x})\| \leq c_{m,n} \|\mathbf{x}\|^{-n/m}\]

where \(\|\cdot\|\) denotes distance to the nearest integer.

\begin{center}\rule{0.5\linewidth}{0.5pt}\end{center}

\section{10. The Roth--Gallai Theorem}\label{ch1958klausfriedrichroth:the-rothgallai-theorem}

\subsection{10.1 Statement}\label{ch1958klausfriedrichroth:statement-1}

The \textbf{Roth--Gallai theorem} (proved with Tibor Gallai) concerns
the existence of arithmetic progressions in subsets of integers. It
states:

\begin{quote}
For any finite coloring of \(\mathbb{Z}\), there exists a monochromatic
three-term arithmetic progression.
\end{quote}

\subsection{10.2 The Color-Focusing
Method}\label{ch1958klausfriedrichroth:the-color-focusing-method}

Roth's proof of the Gallai theorem used a \textbf{color-focusing
argument}: if there is no monochromatic three-term progression, then the
coloring must satisfy strong structural constraints that eventually
force a contradiction. This method was later abstracted in the
\textbf{Hales--Jewett theorem} and its proof.

\subsection{10.3 Generalizations}\label{ch1958klausfriedrichroth:generalizations}

This is a special case of \textbf{van der Waerden's theorem} (1927),
which states that for any \(k\), any finite coloring of \(\mathbb{Z}\)
contains arbitrarily long monochromatic arithmetic progressions. Roth's
Fourier-analytic proof of the three-term case was an important precursor
to the ergodic-theoretic proofs of Szemerédi and Furstenberg.

\subsection{10.4 The Polynomial Van der Waerden
Theorem}\label{ch1958klausfriedrichroth:the-polynomial-van-der-waerden-theorem}

\textbf{Bergelson and Leibman} (1996) proved the polynomial van der
Waerden theorem: for any finite coloring of \(\mathbb{Z}\), there exist
\(x, y\) such that \(x, x+y, x+y^2\) are all the same color. This
generalizes van der Waerden's theorem and uses the structural approach
pioneered by Roth.

\begin{center}\rule{0.5\linewidth}{0.5pt}\end{center}

\section{11. The Davenport--Roth
Estimates}\label{ch1958klausfriedrichroth:the-davenportroth-estimates}

\subsection{11.1 Collaboration with
Davenport}\label{ch1958klausfriedrichroth:collaboration-with-davenport}

Roth's long collaboration with \textbf{Harold Davenport} produced a
series of papers on: - \textbf{Diophantine approximation} and the
geometry of numbers. - \textbf{The theory of quadratic forms} and their
representation of numbers. - \textbf{The distribution of values of
quadratic forms} over integers.

\subsection{11.2 The Davenport--Roth
Theorem}\label{ch1958klausfriedrichroth:the-davenportroth-theorem}

The \textbf{Davenport--Roth theorem} gives estimates for the number of
integer solutions to systems of homogeneous linear equations with
bounded height. These estimates are essential for the
\textbf{Hardy--Littlewood circle method} and its applications.

\subsection{11.3 Davenport--Roth on Quadratic
Forms}\label{ch1958klausfriedrichroth:davenportroth-on-quadratic-forms}

Davenport and Roth also studied the representation of integers by
quadratic forms. Their work on the \textbf{distribution of values of
indefinite quadratic forms} used the geometry of numbers and the theory
of reduction. Key results include: - Bounds for the number of
representations of an integer by a positive definite quadratic form. -
The \textbf{Davenport--Roth inequality} for the number of solutions to
\(Q(x) = m\) with \(|x| \leq N\). - Asymptotic formulas for the number
of integer points on quadrics.

\subsection{11.4 The Hardy--Littlewood Circle
Method}\label{ch1958klausfriedrichroth:the-hardylittlewood-circle-method}

The \textbf{Hardy--Littlewood circle method} is a powerful technique in
analytic number theory for counting solutions to Diophantine equations.
Roth's work on the large sieve and the geometry of numbers provided
crucial estimates for the minor arcs in the circle method.

\begin{center}\rule{0.5\linewidth}{0.5pt}\end{center}

\section{12. Quantitative Roth
Theorems}\label{ch1958klausfriedrichroth:quantitative-roth-theorems}

\subsection{12.1 Effective Roth Theorems}\label{ch1958klausfriedrichroth:effective-roth-theorems}

Over the decades, several authors have improved the quantitative aspects
of Roth's theorem: - \textbf{Bombieri and van der Poorten (1987):} New
proofs using the theory of \textbf{Padé approximation}. -
\textbf{Faltings and Wüstholz (1994):} A new proof using the theory of
\textbf{Mumford compactifications}. - \textbf{Vojta (1996):} A proof
using \textbf{Arakelov theory} and the theory of heights.

\subsection{12.2 The Role of Roth's
Lemma}\label{ch1958klausfriedrichroth:the-role-of-roths-lemma}

\textbf{Roth's lemma} is the key combinatorial component of the proof.
It states: given an integer polynomial \(P\) in \(n\) variables of
degree at most \(d\) and a set of rational points with large common
denominator, if \(P\) vanishes at many of these points, then \(P\) is
reducible in a precise sense. More precisely:

\begin{quote}
\textbf{Roth's Lemma:} Let \(P(x_1, \ldots, x_n)\) be a non-zero
polynomial with integer coefficients and total degree at most \(d\). Let
\(p_1/q, \ldots, p_n/q\) be distinct rational numbers with common
denominator \(q\). If \(P(p_1/q, \ldots, p_n/q) = 0\) for many distinct
\(n\)-tuples, then \(P\) must factor as a product of lower-degree
polynomials.
\end{quote}

This lemma has since found applications in transcendence theory and the
theory of heights.

\subsection{12.3 The Bombieri--van der Poorten
Approach}\label{ch1958klausfriedrichroth:the-bombierivan-der-poorten-approach}

Bombieri and van der Poorten (1987) gave a new proof of Roth's theorem
using \textbf{Padé approximation}. Instead of Roth's combinatorial
lemma, they used the theory of \textbf{rational approximation} to
algebraic functions. This approach connects Roth's theorem to the theory
of \textbf{modular forms} and \textbf{Hecke operators}.

\begin{center}\rule{0.5\linewidth}{0.5pt}\end{center}

\section{13. Roth's Theorem in Transcendence
Theory}\label{ch1958klausfriedrichroth:roths-theorem-in-transcendence-theory}

\subsection{13.1 Connection to
Transcendence}\label{ch1958klausfriedrichroth:connection-to-transcendence}

Roth's theorem is a cornerstone of \textbf{transcendence theory}.
Combined with the theory of linear forms in logarithms (Baker), it
yields transcendence results for large classes of numbers.

\subsection{13.2 The Six Exponentials
Theorem}\label{ch1958klausfriedrichroth:the-six-exponentials-theorem}

The \textbf{six exponentials theorem} (Siegel, Lang, Ramachandra)
states: if \(x_1, x_2\) are linearly independent over \(\mathbb{Q}\) and
\(y_1, y_2, y_3\) are also linearly independent over \(\mathbb{Q}\),
then at least one of the six numbers \(e^{x_i y_j}\) is transcendental.
This result uses Roth-type estimates.

\subsection{13.3 The Schanuel Conjecture}\label{ch1958klausfriedrichroth:the-schanuel-conjecture}

\textbf{Schanuel's conjecture} --- a far-reaching generalization of
Roth's theorem and Baker's theorem --- would imply that: - \(e\) and
\(\pi\) are algebraically independent. - Many numbers of the form
\(\alpha^\beta\) for algebraic \(\alpha, \beta\) are transcendental. -
The Lindemann--Weierstrass theorem generalizes to arbitrary algebraic
independence.

\begin{center}\rule{0.5\linewidth}{0.5pt}\end{center}

\section{14. Roth's Influence on Additive
Combinatorics}\label{ch1958klausfriedrichroth:roths-influence-on-additive-combinatorics}

\subsection{14.1 The Legacy of the Progression-Free
Result}\label{ch1958klausfriedrichroth:the-legacy-of-the-progression-free-result}

Roth's theorem on three-term progressions is the founding result of
\textbf{additive combinatorics}. It inspired: - \textbf{Szemerédi's
theorem (1975):} Sets of positive upper density contain \(k\)-term
arithmetic progressions. - \textbf{Gowers' theorem (2001):} Quantitative
bounds for \(k\)-term progressions using higher-order Fourier analysis.
- \textbf{Green--Tao theorem (2004):} The primes contain arbitrarily
long arithmetic progressions. - \textbf{The polynomial method} in
additive combinatorics (Croot, Lev, Pach; Ellenberg, Gijswijt).

\subsection{14.2 The Density Hales--Jewett
Theorem}\label{ch1958klausfriedrichroth:the-density-halesjewett-theorem}

Roth's Fourier-analytic method was extended to \textbf{higher
dimensions} and \textbf{algebraic structures}: - The \textbf{density
Hales--Jewett theorem} (Furstenberg, Katznelson) is a multidimensional
generalization. - The \textbf{Polymath project} (2009--2010) found a new
combinatorial proof of the density Hales--Jewett theorem.

\subsection{14.3 The Cap Set Problem}\label{ch1958klausfriedrichroth:the-cap-set-problem}

The \textbf{cap set problem} asks for the size of the largest subset of
\(\mathbb{F}_3^n\) containing no three-term arithmetic progression.
Roth's method gave the first non-trivial upper bounds. The problem was
dramatically improved by Croot, Lev, and Pach (2017) using the
\textbf{polynomial method}, leading to bounds of the form \(O(c^n)\) for
some \(c < 3\).

\begin{center}\rule{0.5\linewidth}{0.5pt}\end{center}

\section{15. Roth and the Bombieri--Vinogradov
Theorem}\label{ch1958klausfriedrichroth:roth-and-the-bombierivinogradov-theorem}

\subsection{15.1 The Theorem}\label{ch1958klausfriedrichroth:the-theorem}

The \textbf{Bombieri--Vinogradov theorem} (1965) states that the primes
are well-distributed in arithmetic progressions on average over large
moduli:

\[\sum_{q \leq Q} \max_{a: (a,q)=1} \left|\pi(x; q, a) - \frac{\text{li}(x)}{\phi(q)}\right| = O\left(\frac{x}{(\log x)^A}\right)\]

for \(Q \leq x^{1/2}/(\log x)^{B}\).

\subsection{15.2 Roth's Contribution}\label{ch1958klausfriedrichroth:roths-contribution-1}

Roth's work on the large sieve was a crucial ingredient in Bombieri's
proof. The large sieve inequality provides the key estimate for handling
the error terms when averaging over moduli.

\begin{center}\rule{0.5\linewidth}{0.5pt}\end{center}

\section{16. Roth's Theorem on Integer Solutions of
Equations}\label{ch1958klausfriedrichroth:roths-theorem-on-integer-solutions-of-equations}

\subsection{16.1 Thue Equations}\label{ch1958klausfriedrichroth:thue-equations}

A \textbf{Thue equation} is an equation of the form \(F(x,y) = m\),
where \(F\) is a homogeneous polynomial of degree \(\geq 3\) with
integer coefficients and \(m\) is an integer. Thue (1909) proved that
such equations have finitely many integer solutions.

\subsection{16.2 Roth's Contribution}\label{ch1958klausfriedrichroth:roths-contribution-2}

Roth's theorem provides the key Diophantine approximation estimate used
in the proof of the finiteness of solutions to Thue equations. The type
of estimate needed is:

For an algebraic number \(\alpha\) (a root of \(F(t,1) = 0\)), if
\((x,y)\) is a solution, then \(|x/y - \alpha|\) is small, and Roth's
theorem bounds how small it can be without falling below the
approximation exponent \(2 + \epsilon\).

\subsection{16.3 Effective Solutions}\label{ch1958klausfriedrichroth:effective-solutions}

While Roth's theorem is ineffective, \textbf{Alan Baker's} theory of
linear forms in logarithms provides effective bounds for integer
solutions to Thue equations. The \textbf{Baker--Schmidt theorem} extends
this to more general Diophantine equations.

\begin{center}\rule{0.5\linewidth}{0.5pt}\end{center}

\section{17. Irregularities of
Distribution}\label{ch1958klausfriedrichroth:irregularities-of-distribution}

\subsection{17.1 The Problem}\label{ch1958klausfriedrichroth:the-problem-1}

The \textbf{irregularities of distribution} problem asks: given \(N\)
points in \([0,1]^d\), what is the smallest possible discrepancy? Roth's
1954 paper provided the answer for the \(L^2\) discrepancy.

\subsection{17.2 The Optimality}\label{ch1958klausfriedrichroth:the-optimality}

Roth's lower bound \(\Omega((\log N)^{(d-1)/2}/N)\) is optimal.
Constructions achieving this bound exist: - \textbf{Halton sequences}
(1960) achieve \(O((\log N)^d/N)\). - \textbf{Hammersley sequences} give
the optimal rate. - \textbf{Faure sequences} and \textbf{Sobol
sequences} achieve similar rates.

\subsection{17.3 Applications}\label{ch1958klausfriedrichroth:applications-2}

Discrepancy theory has applications in: - \textbf{Quasi-Monte Carlo
methods} for numerical integration. - \textbf{Computer graphics}
(sampling patterns). - \textbf{Random number generation}. -
\textbf{Computational finance} (option pricing).

\begin{center}\rule{0.5\linewidth}{0.5pt}\end{center}

\section{18. Roth's Work on the ABC
Conjecture}\label{ch1958klausfriedrichroth:roths-work-on-the-abc-conjecture}

\subsection{18.1 The ABC Conjecture}\label{ch1958klausfriedrichroth:the-abc-conjecture}

The \textbf{ABC conjecture} (Masser, Oesterlé) states: for any
\(\epsilon > 0\), there exists \(C_\epsilon\) such that for any coprime
integers \(a + b = c\):

\[\max(|a|, |b|, |c|) \leq C_\epsilon \prod_{p|abc} p^{1+\epsilon}\]

\subsection{18.2 Connection to Roth's
Theorem}\label{ch1958klausfriedrichroth:connection-to-roths-theorem}

Roth's theorem on Diophantine approximation is closely related to the
ABC conjecture. The \textbf{Mason--Stothers theorem} (for polynomials)
is an exact analogue that implies the polynomial ABC conjecture. The ABC
conjecture generalizes Roth's theorem to the setting of three-term
equations.

\subsection{18.3 The Roth Theorem for Number
Fields}\label{ch1958klausfriedrichroth:the-roth-theorem-for-number-fields}

Roth's theorem has been generalized to \textbf{number fields} and
\textbf{function fields}: - \textbf{Lang (1960):} Roth's theorem for
approximation of algebraic numbers by algebraic numbers in a number
field. - \textbf{Uchiyama (1964):} A version for function fields in
characteristic zero. - \textbf{Mahler and Sprindzhuk:} Generalizations
to \(p\)-adic fields.

\begin{center}\rule{0.5\linewidth}{0.5pt}\end{center}

\section{19. The Theory of Heights and Roth's
Theorem}\label{ch1958klausfriedrichroth:the-theory-of-heights-and-roths-theorem}

\subsection{19.1 Heights}\label{ch1958klausfriedrichroth:heights}

The \textbf{height} \(H(\alpha)\) of an algebraic number \(\alpha\)
measures its arithmetic complexity. For a rational number \(p/q\) in
lowest terms, \(H(p/q) = \max(|p|, |q|)\). Roth's theorem can be
restated as: for any algebraic \(\alpha\) and \(\epsilon > 0\):

\[|\alpha - \beta| \leq H(\beta)^{-2-\epsilon}\]

has only finitely many algebraic \(\beta\) of bounded degree.

\subsection{19.2 The Northcott Property}\label{ch1958klausfriedrichroth:the-northcott-property}

\textbf{Northcott's theorem} states that there are only finitely many
algebraic numbers of bounded degree and bounded height. Roth's theorem
refines this by showing that algebraic numbers cannot be too close to a
fixed algebraic number.

\subsection{19.3 The Bogomolov
Conjecture}\label{ch1958klausfriedrichroth:the-bogomolov-conjecture}

The \textbf{Bogomolov conjecture} (proved by Ullmo and Zhang) states
that points of small height on an abelian variety cannot be too dense.
This is a far-reaching generalization of Roth's theorem in the context
of Arakelov geometry.

\begin{center}\rule{0.5\linewidth}{0.5pt}\end{center}

\section{19. Roth's Theorem on Approximation of Algebraic Numbers by
Algebraic
Numbers}\label{ch1958klausfriedrichroth:roths-theorem-on-approximation-of-algebraic-numbers-by-algebraic-numbers}

\subsection{19.1 A Generalization}\label{ch1958klausfriedrichroth:a-generalization}

Roth also proved a generalization: for algebraic numbers \(\alpha\) and
\(\beta\), the inequality:

\[|\alpha - \beta| < H(\beta)^{-2-\epsilon}\]

has only finitely many algebraic \(\beta\) of bounded degree, where
\(H(\beta)\) is the height of \(\beta\).

\subsection{19.2 The Schmidt Subspace
Theorem}\label{ch1958klausfriedrichroth:the-schmidt-subspace-theorem}

\textbf{Wolfgang Schmidt's} subspace theorem generalizes Roth's theorem
to the setting of simultaneous approximation by algebraic numbers. It is
one of the effective results in Diophantine approximation.

\begin{center}\rule{0.5\linewidth}{0.5pt}\end{center}

\section{20. The Roth--Wirsing Theorem}\label{ch1958klausfriedrichroth:the-rothwirsing-theorem}

\subsection{20.1 Statement}\label{ch1958klausfriedrichroth:statement-2}

The \textbf{Roth--Wirsing theorem} (proved with Eduard Wirsing) concerns
the approximation of real numbers by algebraic integers. It states that
for any real \(\alpha\) and any \(\epsilon > 0\), there are infinitely
many algebraic integers \(\beta\) such that:

\[|\alpha - \beta| < H(\beta)^{-2+\epsilon}\]

\subsection{20.2 Optimality}\label{ch1958klausfriedrichroth:optimality}

The exponent \(2\) is optimal, as shown by Roth's theorem on the
approximation of algebraic numbers by rationals. The Roth--Wirsing
theorem shows that algebraic integers are better approximants to real
numbers than algebraic numbers of the same degree.

\begin{center}\rule{0.5\linewidth}{0.5pt}\end{center}

\section{21. Roth's Legacy in Modern Number
Theory}\label{ch1958klausfriedrichroth:roths-legacy-in-modern-number-theory}

\subsection{21.1 The Roth Seminar}\label{ch1958klausfriedrichroth:the-roth-seminar}

Roth's influence extends through his students and collaborators at
University College London and Imperial College. He supervised many
distinguished number theorists, including:

\begin{itemize}
\tightlist
\item
 \textbf{Bryan Birch} (known for Birch--Swinnerton-Dyer conjecture).
\item
 \textbf{Peter Swinnerton-Dyer} (BSD conjecture, sieve methods).
\item
 \textbf{David Masser} (ABC conjecture, transcendence theory).
\item
 \textbf{Robert Vaughan} (analytic number theory, Waring's problem).
\end{itemize}

\subsection{21.2 The Davenport--Roth
School}\label{ch1958klausfriedrichroth:the-davenportroth-school}

The \textbf{Davenport--Roth school} of number theory at UCL and
Cambridge produced a generation of leading number theorists, including
Birch, Swinnerton-Dyer, Vaughan, and Masser. Their collective work
shaped analytic number theory and Diophantine geometry for decades.

\subsection{21.3 Roth's Theorem in the 21st
Century}\label{ch1958klausfriedrichroth:roths-theorem-in-the-21st-century}

Roth's theorem remains an active area of research: - \textbf{Exponential
improvements} in the bounds for progression-free sets (Bloom--Sisask). -
\textbf{New proofs} of Roth's theorem using the polynomial method. -
\textbf{Generalizations} to arbitrary finite fields and non-abelian
groups.

\begin{center}\rule{0.5\linewidth}{0.5pt}\end{center}

\section{22. Roth's Later Work and Final
Years}\label{ch1958klausfriedrichroth:roths-later-work-and-final-years}

\subsection{22.1 Later Papers}\label{ch1958klausfriedrichroth:later-papers}

In his later decades, Roth continued to work on: -
\textbf{Irregularities of distribution} in higher dimensions, refining
his earlier discrepancy bounds. - \textbf{The large sieve} and its
applications to the distribution of primes. - \textbf{Sequences with
restricted growth} and their combinatorial properties.

\subsection{22.2 Retirement and Legacy}\label{ch1958klausfriedrichroth:retirement-and-legacy}

Roth retired from Imperial College in 1988 but remained mathematically
active. He moved to Scotland, where he enjoyed the landscape and quiet
life. He passed away in Inverness on 10 November 2015 at the age of 90.

\subsection{22.3 Quotations}\label{ch1958klausfriedrichroth:quotations}

Carl Ludwig Siegel's comment on Roth's theorem is notable: ``This
result will be remembered as long as mankind is interested in
mathematics.'' Harold Davenport's tribute at the 1958 ICM placed Roth
among the greatest problem-solvers in the history of number theory:
``The great unsolved problems of mathematics may still yield to direct
attack, however difficult and forbidding they appear to be.''

\begin{center}\rule{0.5\linewidth}{0.5pt}\end{center}

\section{23. Awards and Recognition}\label{ch1958klausfriedrichroth:awards-and-recognition}

\textbf{1958:} - Fields Medal (International Mathematical Union) ---
awarded at the ICM in Edinburgh; first British recipient.

\textbf{1960:} - Elected Fellow of the Royal Society (FRS).

\textbf{1979:} - Elected Fellow of University College London.

\textbf{1983:} - De Morgan Medal (London Mathematical Society).

\textbf{1988:} - Knighted? (No --- Roth refused knighthood.)

\textbf{1989:} - Honorary Fellow, Peterhouse, Cambridge.

\textbf{1991:} - Sylvester Medal (Royal Society) --- ``For his many
contributions to number theory and in particular his solution of the
famous problem concerning approximating algebraic numbers by
rationals.''

\textbf{1993:} - Honorary Fellow of the Royal Society of Edinburgh.

\textbf{1999:} - Fellow of Imperial College London.

\textbf{Elected to:} - Royal Society (London). - Royal Society of
Edinburgh (Honorary Fellow). - London Mathematical Society (member 1951,
De Morgan Medal 1983). - Honorary Fellow, Peterhouse, Cambridge. -
Fellow of University College London. - Fellow of Imperial College
London.

\begin{center}\rule{0.5\linewidth}{0.5pt}\end{center}

\section{23. Selected Major
Publications}\label{ch1958klausfriedrichroth:selected-major-publications}

\begin{enumerate}
\def\labelenumi{\arabic{enumi}.}
\item
 \textbf{``On certain sets of integers.''} \emph{Journal of the London
 Mathematical Society} 28, 104--109, 1953. --- \emph{Roth's theorem on
 sets with no three-term arithmetic progression.}
\item
 \textbf{``On the discrepancy of sequences in the unit interval.''}
 \emph{Journal of the London Mathematical Society} 29, 378--385, 1954.
 --- \emph{Roth's optimal lower bound for discrepancy.}
\item
 \textbf{``Rational approximations to algebraic numbers.''}
 \emph{Mathematika} 2, 1--20, 1955 (Corrigendum, \emph{Mathematika} 2,
 168, 1955). --- \emph{Roth's theorem on Diophantine approximation.}
\item
 \textbf{``On the large sieves of Linnik and Rényi.''} \emph{Journal of
 the London Mathematical Society} 39, 309--320, 1964. --- \emph{Seminal
 contribution to the large sieve method.}
\item
 \textbf{``The distribution of integers having no small prime
 factors.''} \emph{Journal of the London Mathematical Society} 39,
 321--328, 1964. --- \emph{Refined sieve estimates.}
\item
 \textbf{``A problem in combinatorial number theory.''} \emph{Acta
 Arithmetica} 12, 167--176, 1966. --- \emph{Further results on
 arithmetic progressions in integer sets.}
\item
 \textbf{``Irregularities of distribution.''} \emph{Mathematika} 19,
 178--186, 1972. --- \emph{Higher-dimensional discrepancy bounds.}
\item
 \textbf{``Rational approximations to algebraic numbers''} (with D.
 Ridout). \emph{Mathematika} 4, 125--131, 1957. --- \emph{The
 Roth--Ridout theorem.}
\item
 \textbf{``The large sieve and its applications''} (with H. Davenport).
 \emph{Mathematika} 12, 193--204, 1965. --- \emph{Davenport--Roth large
 sieve results.}
\item
 \textbf{``A problem on irregular distributions.''} \emph{Colloquium
 Mathematicum} 14, 290--294, 1966. --- \emph{Further discrepancy
 results.}
\end{enumerate}

\begin{center}\rule{0.5\linewidth}{0.5pt}\end{center}

\begin{center}\rule{0.5\linewidth}{0.5pt}\end{center}

\section*{Exercises}
\addcontentsline{toc}{section}{Exercises}

\begin{exercise} Prove Liouville's theorem: if $\alpha$ is irrational and algebraic of degree $n$, then $|\alpha - p/q| > c/q^n$ for all $p/q$. \end{exercise}

\begin{exercise} Show that Roth's theorem implies that algebraic irrationals have irrationality measure 2 (best possible). \end{exercise}

\begin{exercise} Verify that $\sum_{q=1}^{Q} \|q\alpha\| \leq C \log Q$ for any irrational $\alpha$ using Dirichlet's approximation. \end{exercise}

\begin{exercise} Compute the continued fraction expansion of $\sqrt{2}$ and discuss how it relates to best rational approximations. \end{exercise}

\begin{center}\rule{0.5\linewidth}{0.5pt}\end{center}

\section*{Summary}
\addcontentsline{toc}{section}{Summary}

Roth proved that for any algebraic irrational $\alpha$ and any $\epsilon > 0$, the inequality $|\alpha - p/q| < 1/q^{2+\epsilon}$ has only finitely many solutions. This settled the Thue-Siegel conjecture and remains the strongest result in metric diophantine approximation.

\section{24. References and Further
Reading}\label{ch1958klausfriedrichroth:references-and-further-reading}

\subsection{Official Sources}\label{ch1958klausfriedrichroth:official-sources}

\begin{itemize}
\tightlist
\item
 IMU Fields Medal 1958 Announcement:
 https://www.mathunion.org/imu-awards/fields-medal/fields-medals-1958
\item
 Royal Society Biographical Memoir (Chen and Vaughan, 2017).
\item
 LMS Biography: Klaus Friedrich Roth, 1925--2015 (Bulletin of the LMS,
 2018).
\end{itemize}

\subsection{Books}\label{ch1958klausfriedrichroth:books}

\begin{itemize}
\tightlist
\item
 Roth, K.F. \emph{Rational Approximations to Algebraic Numbers}. LMS
 Lecture Note Series, 1955.
\item
 Roth, K.F. \emph{Collected Works}. Oxford University Press, 1998.
\item
 Davenport, H. \emph{Multiplicative Number Theory}. Springer, 1980.
\item
 Hindry, M. and Silverman, J. \emph{Diophantine Geometry: An
 Introduction}. Springer, 2000.
\item
 Lang, S. \emph{Diophantine Approximation}. Springer, 1983.
\item
 Schmidt, W. \emph{Diophantine Approximation and Diophantine
 Equations}. Springer, 1991.
\end{itemize}

\subsection{Key Articles}\label{ch1958klausfriedrichroth:key-articles}

\begin{itemize}
\tightlist
\item
 Davenport, H. ``The Work of Klaus Roth.'' ICM 1958 Laudatio.
\item
 Siegel, C.L. ``Über die Approximation algebraischer Zahlen.'' (Comment
 on Roth's theorem, 1956.)
\item
 Bombieri, E. and van der Poorten, A.J. ``Roth's theorem and the theory
 of heights.'' \emph{Journal of Number Theory} 24, 1986.
\item
 Chen, W. and Vaughan, R. ``Klaus Friedrich Roth, 29 October 1925 -- 10
 November 2015.'' \emph{Biographical Memoirs of Fellows of the Royal
 Society} 64, 91--109, 2018.
\item
 Chen, W. ``The work of Klaus Roth on irregularities of distribution.''
 \emph{Functional Analysis and Its Applications} 34, 2000.
\end{itemize}

\subsection{Textbooks and Surveys}\label{ch1958klausfriedrichroth:textbooks-and-surveys}

\begin{itemize}
\tightlist
\item
 Tao, T. and Vu, V. \emph{Additive Combinatorics}. Cambridge University
 Press, 2006.
\item
 Kuipers, L. and Niederreiter, H. \emph{Uniform Distribution of
 Sequences}. Wiley, 1974.
\item
 Schmidt, W. \emph{Diophantine Approximation}. Springer, 1980.
\item
 Gowers, T. ``A new proof of Szemerédi's theorem.'' \emph{GAFA} 11,
 2001.
\item
 Bloom, T. and Sisask, O. ``Breaking the logarithmic barrier in Roth's
 theorem.'' \emph{Annals of Mathematics} 192, 2020.
\end{itemize}

\subsection{Further Reading}\label{ch1958klausfriedrichroth:further-reading}

\begin{itemize}
\tightlist
\item
 Matoušek, J. \emph{Geometric Discrepancy}. Springer, 1999.
\item
 Chazelle, B. \emph{The Discrepancy Method}. Cambridge University
 Press, 2000.
\item
 Graham, R. and Nesetril, J. \emph{The Mathematics of Paul Erdős}.
 Springer, 1997.
\end{itemize}

\subsection{Biography}\label{ch1958klausfriedrichroth:biography}

\begin{itemize}
\tightlist
\item
 Chen, W. and Vaughan, R. ``Klaus Friedrich Roth, 29 October 1925 -- 10
 November 2015.'' \emph{Biographical Memoirs of Fellows of the Royal
 Society} 64, 91--109, 2018.
\item
 ``Klaus Friedrich Roth 1925--2015.'' \emph{Bulletin of the London
 Mathematical Society} 50, 729--739, 2018.
\end{itemize}

\begin{center}\rule{0.5\linewidth}{0.5pt}\end{center}


\begin{center}\rule{0.5\linewidth}{0.5pt}\end{center}

\section*{See Also}
\addcontentsline{toc}{section}{See Also}

This chapter relates to the following work by other Fields Medalists:
\begin{itemize}
  \item \textbf{\hyperref[ch1986gerdfaltings:gerd-faltings-fields-medal-1986-a-comprehensive-report]{See Chapter on related work}} --- Faltings extended Roth's theorem.
  \item \textbf{\hyperref[ch1970alanbaker:alan-baker-fields-medal-1970-a-comprehensive-report]{See Chapter on related work}} --- Baker's work on linear forms relates to Roth.
\end{itemize}

\emph{(End of report)}
